\begin{proof}
(a) Composite cycles. By Proposition \ref{prop:c10-compbound} the composite cycles through $(A|_{W_{\mathbb R}^\perp},A|_{W_{\mathbb R}^\perp})$ whose image under $Q$ has the two outputs $A|_{W_{\mathbb R}^\perp}$ are the two triangles formed by the vertex $A|_V(A|_{W_{\mathbb R}^\perp})^2$, entered at $A|_{W_{\mathbb R}^\perp}$, left at $A|_V$, with the output $A|_{W_{\mathbb R}^\perp}$, and the vertex $([A|_{W_{\mathbb R}^\perp},c],A^*|_{W_{\mathbb R}^\perp})$, entered at $c$, left at $A|_{W_{\mathbb R}^\perp}$, with the output $A^*|_{W_{\mathbb R}^\perp}$, joined by the composite line in either order; a cycle with two outputs has two cubic vertices and the composite line between them, since the line ends on no letter $A|_{W_{\mathbb R}^\perp}$. Both vertices act as the identity on the $\mathfrak{so}(W_{\mathbb R}^\perp)$ index along the cycle, so each triangle maps $(\iota_{w_1}A,\iota_{w_2}A)$ to outputs with the fibre elements $w_1$, $w_2$ in the order of the letters. The leg sets $A|_V\,A|_{W_{\mathbb R}^\perp}A^*|_{W_{\mathbb R}^\perp}$ and $A|_{W_{\mathbb R}^\perp}^2A^*|_V$ give outputs with a factor $A|_V$, the covariant completion of $x\cdot\partial$, outside this part.

(b) Form. A local $\mathrm{Spin}(4)$-invariant density of weight $2$, the position at weight $-1$, at most linear in the position, with the outputs $\iota_{w_1}A\,\iota_{w_2}A$ and no further derivatives, is $b_1\,\iota_{w_1}A\,\iota_{w_2}A+b_2\,(x\cdot\partial\,\iota_{w_1}A)\,\iota_{w_2}A+b_3\,\iota_{w_1}A\,(x\cdot\partial\,\iota_{w_2}A)$ in the planar adjacent form of Lemma \ref{lem:c10-colour} (iv), with numbers $b_1,b_2,b_3$.

(c) Normalisation. For the bilinear $\int\omega\,\operatorname{tr}(\iota_{w_1}A\,\iota_{w_2}A)$ the single pair is adjacent on both sides: by Lemma \ref{lem:c10-colour} (ii) with $M_1=1$ and $r=2$ the terms $S=\emptyset$ and $S=\{1,2\}$ give $N\operatorname{tr}(Y_1Y_2)$ and $N\operatorname{tr}(Y_2Y_1)$, and the others vanish for traceless outputs, so the bilinear receives the density on $(X_1,X_2)$ and on $(X_2,X_1)$: $2b_1O+(b_2+b_3)\sum_lx\cdot\partial X_l$. Proposition \ref{prop:c9-triangle} gives $-2$ units times $v.O=\sum_l(2x\cdot\partial+2)X_l$ for the traceless and the trace parts alike (the vertices act as the identity on $\mathfrak{so}(W_{\mathbb R}^\perp)$, table of Theorem \ref{thm:c9-D1}): $b_1=-4$ and $b_2+b_3=-4$ units.

(d) Length $L$. Summing over the $L$ cyclically adjacent pairs gives $b_1L\,O+(b_2+b_3)\sum_lx\cdot\partial X_l=-2\,v.O$ units; the part $\tfrac{b_2-b_3}2\big[(x\cdot\partial X_a)X_b-X_a(x\cdot\partial X_b)\big]$ sums to zero, each letter receiving $x\cdot\partial$ once with each sign. So the composite term of the pair is $-v.(X_aX_b)$ units modulo such terms.

(e) One body. By Theorem \ref{thm:c10-D1} each letter receives $(\tfrac53-\theta)v.X$ units beyond $-8X$; with (d), the composite term of $O$ minus the bracket with $\hbar F_1^{(1)}(\Delta)$ minus $4$ units times the field-counting functional is $(\tfrac53-\theta-2)v.O$ units. For the density, $4$ times the map of Theorem \ref{thm:c10-brst} per pair is $-8$ units times the bracket of Proposition \ref{prop:c10-so6}, and $v.O$ is half the sum over the adjacent pairs of $v.(X_aX_b)$.

(f) Translations. By Proposition \ref{prop:c10-P1} and Theorem \ref{thm:c9-P1}, the composite line along $P_b$ carries no position, the density is $\nu_2(\partial_bX_a)X_b+\nu_3X_a\partial_bX_b$ with numbers $\nu_2$, $\nu_3$, the bilinear (the composite triangle along $P_b$ summed over the two orderings, $-2$ units times $\{S^{(1)}(P_b),O\}$) gives $\nu_2+\nu_3=-2$ units, and the one-body part is $(\tfrac53-\theta_P)\partial_bX$ units; (d) gives $P_1=-(\tfrac13+\theta_P)$ units times $\partial_bO$. (iii): Theorem \ref{thm:c9-theta2}, Corollary \ref{cor:c10-beisert} and Proposition \ref{prop:c9-scheme}, the change along translations being $\tfrac a2\lambda_{P_b}\{S^{(1)}(P_b),-\}$ with $\lambda_{P_b}=-2$, and zero along $\Delta$.
\end{proof}
